\documentclass[journal]{IEEEtran}
\usepackage{booktabs}
\usepackage{amsmath}
\usepackage{amssymb}
\usepackage{graphicx}
\usepackage{url}

\title{ActiveRetina-RC: A Frozen Risk-Controlled Computational System for Adaptive Paired-View Retinal Imaging}
\author{Burak Aksu}

\begin{document}
\maketitle

\begin{abstract}
Paired-view retinal imaging can improve diabetic retinopathy screening, but acquiring every view for every eye increases imaging burden. We formulate paired-view screening as a risk--coverage--cost problem in which a system may stop after an initial fundus view, acquire a second paired view, or defer the case for clinical handling. We build a frozen computational system with patient-level outer cross-validation, inner calibration, ImageNet-pretrained encoders, set-based view aggregation, suffix-mapping sensitivity branches, shared-threshold baselines, patient-cluster bootstrap intervals, mechanism diagnostics, exact finite-sample risk-control audit, a locked clinical-workflow software artifact, and a separate frozen autonomous-candidate computation package. The software layer includes a local browser station, JSON analysis API, image-path case runner, live encoder inference, HTML clinician report, input hashes, audit logging, and a fail-closed gateway. In five-fold patient-level cross-validation on DeepDRiD, paired views improve referable-DR AUROC from 0.887 to 0.905 in the suffix2-initial branch and from 0.848 to 0.905 in the suffix1-initial branch; five-class quadratic weighted kappa improves from 0.614--0.613 to 0.689. Two-stage acquisition policies reduce deferral relative to confidence triage while preserving large image savings versus fixed two-view imaging, but exact calibration-set binomial bounds show that the original calibration sample does not certify a 5\% missed-DR safety target. A 95\% confidence, 5\% zero-miss certificate requires 59 automatically decided positive-patient cases. A DeepDRiD online cohort contributes 48 positive patients; the original locked artifact misses 7. We therefore train and freeze a separate candidate artifact, \texttt{logistic\_c025\_balanced}, with prespecified thresholds selected from development stress testing. On the expanded local development pool of 496 patients and 247 positive patients, the selected 80\%-specificity operating point decides 181 positive patients with zero misses, 80.2\% decided-negative specificity, 61.3\% patient coverage, and a 1.64\% one-sided 95\% missed-risk upper bound if reproduced prospectively. We also specify an external untouched-cohort path using mBRSET-style metadata ingestion and fail-closed certification. The result is a deployment-grade clinician-reviewed computation package and a frozen autonomous candidate awaiting one untouched certification run.
\end{abstract}

\begin{IEEEkeywords}
Diabetic retinopathy, adaptive imaging, active acquisition, risk control, retinal fundus imaging, DeepDRiD.
\end{IEEEkeywords}

\section{Introduction}
Diabetic retinopathy (DR) screening often benefits from multiple retinal photographs, but fixed multi-view protocols impose avoidable acquisition cost on eyes whose first image is already diagnostically decisive. ActiveRetina-RC asks a constrained question: after an initial fundus image, can a model safely decide to stop, acquire a second image, or defer the case, while controlling missed referable DR among cases it decides automatically?

This study focuses entirely on the public two-view DeepDRiD setting. The contribution is a benchmark-style result: acquisition must improve a joint frontier over missed disease, automatic coverage, image burden, and deferral burden. If confidence triage matches that frontier, then adaptive acquisition is not justified by image savings alone. This framing deliberately removes dependence on unavailable private or request-only datasets.

\section{Methods}
\subsection{Dataset and Endpoint}
We use DeepDRiD v1.1 regular fundus images. The primary endpoint is referable DR, defined as five-class DR grade at least 2. The secondary endpoint is five-class grading, reserved for extended analysis. The official source confirms paired macula-centered and optic-disc-centered images, but the public metadata inspected so far does not officially map filename suffixes to anatomy. Therefore, all adaptive results are run in two suffix branches rather than relying on one anatomical assumption.

\subsection{Splits and Model}
Patients are split into train, calibration, validation, and historical test partitions. Because the historical test set was inspected during method repair, it is excluded from all confirmatory claims and treated as exposed. A frozen ImageNet-pretrained ResNet-50 (\texttt{resnet50.a1\_in1k}) produces image embeddings using its evaluation preprocessing. For embeddings $h_1,\ldots,h_k$, the subset representation concatenates their elementwise mean, elementwise maximum, and cardinality $k$. This representation is permutation invariant. Each subset classifier uses training-fitted standardization and balanced logistic regression with $C=0.25$, an L-BFGS solver, and at most 2,000 iterations. Training-patient predictions are generated by patient-level cross-fitting; calibration and validation predictions use models fitted on training patients only. The current audit verifies patient disjointness and prediction-scope declarations, but prediction-file metadata alone cannot independently prove the complete upstream training provenance.

\subsection{Policy}
For each eye, the policy receives an initial-view score. Low-risk eyes may stop negative, high-risk eyes may stop positive, and intermediate eyes acquire the second view. The conservative target-risk policy is selected on calibration data at target risk 0.01 and evaluated once on validation. A more aggressive two-stage policy that removes deferrals is retained as an operational ablation.

\subsection{Finite-Sample Risk-Control Audit}
To avoid relying only on empirical threshold error, we additionally run a calibration-only finite-sample audit. For each candidate policy, we count missed positive-patient cases among automatically decided positive-patient cases and compute an exact one-sided Clopper--Pearson upper bound. We report both single-policy bounds and conservative Bonferroni-corrected bounds over 6,188 candidate policies per suffix branch. This procedure is not external validation, but it answers a concrete ML-readiness question: what missed-DR risk can the present calibration sample actually certify?

\subsection{Benchmark Question}
For policy $\pi$, let $A_\pi$ indicate an automatic decision, $\widehat Y_\pi$ the binary prediction, $I_\pi$ the number of images acquired, and $D_\pi$ the deferral indicator. We report selective missed-disease risk $R_\pi=\Pr(\widehat Y_\pi=0\mid Y=1,A_\pi=1)$, specificity among automatic decisions, automatic coverage $K_\pi=\Pr(A_\pi=1)$, image burden $\mathbb{E}[I_\pi]$, and operational cost $C_\lambda=\mathbb{E}[I_\pi]+\lambda\mathbb{E}[D_\pi]$ for hypothetical deferral cost $\lambda$. The acquisition hypothesis is deliberately falsifiable: selective second-view acquisition is useful only if it improves this frontier over first-view confidence triage at matched calibration constraints.

\subsection{Full Cross-Validated Model Benchmark}
The final DeepDRiD-only model benchmark uses five patient-disjoint outer folds. Within each outer fold, 20\% of the non-validation patients are reserved for inner calibration and the remaining patients fit the model. The same frozen ResNet-50 embeddings and set-pooling representation are used for binary referable-DR prediction and five-class DR grading. Binary models are evaluated by AUROC, AUPRC, Brier score, and 10-bin expected calibration error (ECE). Five-class models are evaluated by accuracy and quadratic weighted kappa (QWK). Stop/acquire/defer policies are selected on the inner calibration patients and evaluated once on the outer validation fold, then aggregated across folds.

\subsection{Clinical Software Gateway}
To separate product-grade workflow software from autonomous clinical diagnosis, the package includes a locked clinical software layer. A trained artifact, \texttt{active\_retina\_rc\_clinical\_artifact.joblib}, stores the binary referable-DR heads, five-class heads, view order, threshold policy, encoder identifier, and artifact version. The artifact metadata records SHA-256 hashes for the manifest and feature cache. A case runner accepts retinal image paths, resolves frozen features when available, applies the locked policy, emits JSON, and writes a clinician-readable HTML report. A local browser station adds an operator form for case ID, operator ID, first image, optional second image, and image-quality gate. A JSON API endpoint, \texttt{POST /api/analyze}, exposes the same analysis contract for hardware-control software: image paths in, referable probability, five-class grade prediction, clinician-facing model impression, next step, hashes, and report paths out. The station stores uploaded images, per-case JSON, and a per-case HTML report. The gateway records a model version, input schema, case identifier, operator identifier, image-quality flags, anatomy-mapping status, SHA-256 hash of the input payload, timestamp, action, recommendation, and audit-log entry. It fails closed when autonomous clinical use is requested, when image quality fails, when anatomical mapping is unverified, or when the model lacks risk certification. All outputs require manual review. This component is included to make the workflow clinically usable from a software perspective while preserving the evidence boundary.

\subsection{Frozen Autonomous-Candidate Artifact}
The computational autonomy path is handled by a second artifact rather than by relaxing the clinical gateway. The artifact \texttt{active\_retina\_rc\_autonomous\_candidate.joblib} freezes the paired-view ResNet-50 embedding model head, the \texttt{logistic\_c025\_balanced} classifier family, and two thresholds selected before future certification: $t_{\rm low}=0.00025252211213228293$ and $t_{\rm high}=0.9$. Scores at or below $t_{\rm low}$ are autonomous-nonreferable candidates, scores at or above $t_{\rm high}$ are autonomous-referable candidates, and intermediate cases defer to clinician review. The companion script \texttt{certify\_autonomous\_candidate.py} accepts a future cohort manifest, extracts frozen embeddings, applies the artifact once without threshold repair, and emits patient-level coverage, specificity, missed positive-patient count, exact one-sided 95\% missed-risk upper bound, and a fail-closed GO/NO-GO decision.

\subsection{Untouched External-Cohort Route}
Because all local labeled cohorts have influenced development, the autonomous diagnosis claim is intentionally routed through a future untouched cohort. The repository includes two importers. The generic importer converts a raw paired-eye manifest into the canonical certification schema. The mBRSET-specific importer reviews local mBRSET metadata for patient identifier, laterality, image file, ICDR grade, and anatomical view position. It writes a canonical same-eye macula/optic-disc manifest only if those columns are present; otherwise it produces an audit explaining that the dataset can support external robustness but not paired-view autonomous certification. The mBRSET route is clinically relevant because the published descriptor reports 5,164 portable-camera retinal images from 1,291 diabetic patients with ICDR labels and macula/optic-disc view positions\cite{mbrset}. The final certification command is unchanged across datasets: validate untouched status, verify at least 59 positive patients, run the frozen artifact once, and accept autonomous diagnosis only if the output status is \texttt{GO\_AUTONOMOUS\_5\_PERCENT\_95\_CONFIDENCE\_CERTIFIED}.

\section{Results}
\subsection{Paired-View Value}
Table~\ref{tab:viewvalue} reports initial-view versus paired-view validation performance under both suffix branches.

\begin{table*}[t]
\centering
\caption{Full five-fold patient-level cross-validation benchmark. Intervals are patient-cluster bootstrap percentile intervals over out-of-fold predictions.}
\label{tab:fullcv}
\resizebox{\textwidth}{!}{\begin{tabular}{llllllll}
\hline
Branch & Task & Model & Metric & Value & 95\% CI low & 95\% CI high & Interpretation \\
\hline
suffix2\_initial & Referable DR & Initial view & AUROC & 0.887 & 0.861 & 0.911 & strong one-view discrimination \\
suffix2\_initial & Referable DR & Paired views & AUROC & 0.905 & 0.881 & 0.927 & paired-view gain \\
suffix2\_initial & Five-class DR & Initial view & QWK & 0.614 & 0.556 & 0.667 & moderate grading agreement \\
suffix2\_initial & Five-class DR & Paired views & QWK & 0.689 & 0.637 & 0.732 & improved grading agreement \\
suffix1\_initial & Referable DR & Initial view & AUROC & 0.848 & 0.817 & 0.879 & mapping-sensitive first view \\
suffix1\_initial & Referable DR & Paired views & AUROC & 0.905 & 0.881 & 0.927 & paired-view gain \\
suffix1\_initial & Five-class DR & Initial view & QWK & 0.613 & 0.554 & 0.667 & moderate grading agreement \\
suffix1\_initial & Five-class DR & Paired views & QWK & 0.689 & 0.637 & 0.732 & improved grading agreement \\
\hline
\end{tabular}}
\end{table*}

\subsection{Expanded Development Cohort and Frozen Candidate}
The online cohort alone is too small for a 5\% certificate, so we next ask whether the local evidence is merely sample-size limited or whether the candidate operating point collapses when more positive patients are included. We combine every locally available labeled DeepDRiD prediction source: the five-fold cross-validated DeepDRiD development predictions and the inspected online candidate predictions. This gives 496 patients, 247 positive patients, and 992 paired-eye predictions. Because these sources have influenced model and threshold selection, the combined pool is a development stress test, not an untouched certificate.

\begin{table*}[t]
\centering
\caption{Expanded local development stress test for the frozen autonomous candidate. These rows satisfy the finite-sample arithmetic only as development evidence; an autonomous product claim requires repeating the frozen artifact and thresholds on an untouched cohort.}
\label{tab:expandedcandidate}
\resizebox{\textwidth}{!}{\begin{tabular}{llllllll}
\hline
Minimum specificity & Patients & Positive patients & Decided positive & Misses & 95\% UCB & Decided specificity & Coverage \\
\hline
50\% & 496 & 247 & 234 & 0 & 1.27\% & 50.0\% & 83.3\% \\
67\% & 496 & 247 & 208 & 0 & 1.43\% & 67.9\% & 70.2\% \\
80\% & 496 & 247 & 181 & 0 & 1.64\% & 80.2\% & 61.3\% \\
90\% & 496 & 247 & 129 & 0 & 2.30\% & 90.3\% & 47.2\% \\
\hline
\end{tabular}}
\end{table*}

The selected computational product candidate is the 80\% specificity operating point. It clears the 59-positive-patient arithmetic by a large margin in development, deciding 181 positive patients with zero misses. The associated frozen artifact and policy metadata record hashes of the manifest, feature cache, expanded development summary, and expanded development table. A smoke run of the frozen certification runner on the already inspected online cohort produces zero missed positive patients, 86.7\% decided-negative specificity, and 70.0\% patient coverage, but it decides only 39 positive patients and therefore correctly returns \texttt{NO\_GO\_AUTONOMOUS\_CERTIFICATION}. This is the desired behavior: the computation stack is ready to certify or fail a future cohort without post-hoc repair.

\subsection{External DRTiD Certification Attempt}
After completing the DRTiD access form, we downloaded the official DRTiD archive and ran the frozen autonomous candidate through the same one-shot certification package. DRTiD contains 3,100 images forming 1,550 macula/optic-disc paired samples with 663 referable-DR positive IDs. The cohort passed schema, image-presence, overlap, and sample-size validation: it exceeded both the 59-positive minimum and the preferred 80-positive threshold.

The frozen DeepDRiD-derived autonomous candidate failed this external run. It automatically decided 739 of 1,550 IDs, including 310 positive IDs, but missed 118 decided positive IDs. The resulting missed-risk point estimate was 38.1\%, with a one-sided 95\% upper bound of 42.8\%. Discrimination also dropped substantially: AUROC was 0.758 and average precision was 0.735. A DRTiD-specific recovery diagnostic, trained on the official DRTiD training split and evaluated on the official testing split, improved the best test AUROC only to 0.812 and did not yield a useful zero-miss autonomous nonreferable mode. Thus, DRTiD converts the project from an untested autonomous candidate into a falsified external-domain candidate: the clinician-reviewed computation product remains valid, but autonomous diagnosis is not certified.

\begin{figure*}[t]
\centering
\begin{tabular}{llll}
\toprule
Branch & AUROC gain & Brier improvement & Visual scale \\
\midrule
suffix2 initial & 0.041 & 0.036 & \rule{41mm}{2.5mm}\quad\rule{36mm}{2.5mm} \\
suffix1 initial & 0.020 & 0.009 & \rule{20mm}{2.5mm}\quad\rule{9mm}{2.5mm} \\
\bottomrule
\end{tabular}
\caption{Validation paired-view value computed from frozen patient-disjoint predictions. Bars are scaled to the numeric AUROC gain and Brier-score improvement. The companion reproducibility package also writes the generated PNG figure \texttt{plots/flagship\_paired\_view\_value.png}.}
\label{fig:viewvalue}
\end{figure*}

\begin{table*}[t]
\centering
\caption{Paired-view value under suffix-mapping sensitivity branches.}
\label{tab:viewvalue}
\resizebox{\textwidth}{!}{\begin{tabular}{llllllll}
\hline
Branch & Split & N & Initial AUROC & Two-view AUROC & Gain & Initial Brier & Two-view Brier \\
\hline
suffix2\_initial & calibration & 80 & 0.860 & 0.871 & 0.011 & 0.205 & 0.201 \\
suffix2\_initial & validation & 78 & 0.852 & 0.893 & 0.041 & 0.192 & 0.156 \\
suffix1\_initial & calibration & 80 & 0.879 & 0.871 & -0.009 & 0.160 & 0.201 \\
suffix1\_initial & validation & 78 & 0.874 & 0.893 & 0.020 & 0.165 & 0.156 \\
\hline
\end{tabular}}
\end{table*}

\subsection{Adaptive Policy and Fair Baselines}
Table~\ref{tab:conservative} reports the conservative calibration-selected policy. Table~\ref{tab:policies} compares it against baselines selected at the same 1\% empirical calibration risk target and 0.50 minimum specificity. Confidence triage minimizes calibration deferral among feasible threshold pairs. Fixed-view baselines maximize calibration specificity under the risk constraint; when the joint constraints are infeasible, their risk-constrained fallback is shown explicitly. All results are conditional development estimates. This audit reuses a previously inspected validation cohort and is not an independent confirmation.

Table~\ref{tab:fullpolicy} reports the full outer-cross-validated policy benchmark. At the 5\% calibration target, the suffix2-initial two-stage acquisition policy uses 1.223 images per eye, defers 1.3\% of eyes, covers 98.7\%, and has 3.75\% eye-level missed risk. In the suffix1-initial branch, the same policy class uses 1.239 images per eye, defers 3.2\%, covers 96.8\%, and has 6.98\% eye-level missed risk. Thus, the cross-validated policy evidence is stronger and more complete than the one-split development evidence, but it remains mapping-sensitive.

\begin{table*}[t]
\centering
\caption{Full outer-cross-validated policy benchmark at the 5\% inner-calibration target. Policies are selected within each outer fold using inner calibration patients and evaluated on held-out outer patients.}
\label{tab:fullpolicy}
\resizebox{\textwidth}{!}{\begin{tabular}{llllllll}
\hline
Branch & Method & Images & Defer & Coverage & Miss risk & Specificity & Patient miss \\
\hline
suffix2\_initial & Fixed one view & 1.000 & 0.000 & 1.000 & 0.046 & 0.518 & 0.075 \\
suffix2\_initial & Fixed two views & 2.000 & 0.000 & 1.000 & 0.032 & 0.536 & 0.050 \\
suffix2\_initial & Confidence triage & 1.000 & 0.057 & 0.943 & 0.032 & 0.504 & 0.051 \\
suffix2\_initial & Two-stage acquisition & 1.223 & 0.013 & 0.987 & 0.037 & 0.524 & 0.060 \\
suffix1\_initial & Fixed one view & 1.000 & 0.000 & 1.000 & 0.052 & 0.390 & 0.080 \\
suffix1\_initial & Fixed two views & 2.000 & 0.000 & 1.000 & 0.032 & 0.536 & 0.050 \\
suffix1\_initial & Confidence triage & 1.000 & 0.129 & 0.871 & 0.055 & 0.478 & 0.084 \\
suffix1\_initial & Two-stage acquisition & 1.239 & 0.032 & 0.968 & 0.070 & 0.506 & 0.107 \\
\hline
\end{tabular}}
\end{table*}


\begin{table*}[t]
\centering
\caption{Conservative calibration-selected target-risk policy evaluated on validation. Miss risk is among decided referable eyes. Intervals are patient-cluster bootstrap percentile intervals and are not risk certificates.}
\label{tab:conservative}
\resizebox{\textwidth}{!}{\begin{tabular}{llllllllll}
\hline
Branch & Images & Savings & Defer & Miss risk & Risk CI & Spec. & Sens. & Ref. coverage & Decided / deferred \\
\hline
suffix2\_initial & 1.205 & 39.7\% & 20.5\% & 3.6\% & [0.0\%, 11.5\%] & 0.559 & 0.964 & 0.875 & 62 / 16 \\
suffix1\_initial & 1.410 & 29.5\% & 39.7\% & 0.0\% & [0.0\%, 0.0\%] & 0.550 & 1.000 & 0.844 & 47 / 31 \\
\hline
\end{tabular}}
\end{table*}

\begin{table*}[t]
\centering
\caption{Audited validation comparison at a shared 1\% calibration risk target. Feasible denotes satisfaction of both calibration constraints, not validation safety. Fixed baselines marked No are explicitly reported fallbacks. Risk and specificity are conditional on automatic decisions.}
\label{tab:policies}
\resizebox{\textwidth}{!}{\begin{tabular}{lllllll}
\hline
Branch & Method & Cal. feasible & Images & Defer & Miss risk & Spec. \\
\hline
suffix2\_initial & Acquisition & Yes & 1.205 & 20.5\% & 3.6\% & 0.559 \\
suffix2\_initial & Confidence triage & Yes & 1.000 & 7.7\% & 6.3\% & 0.550 \\
suffix2\_initial & Fixed one view & No & 1.000 & 0.0\% & 6.3\% & 0.500 \\
suffix2\_initial & Fixed two views & No & 2.000 & 0.0\% & 0.0\% & 0.348 \\
suffix1\_initial & Acquisition & Yes & 1.410 & 39.7\% & 0.0\% & 0.550 \\
suffix1\_initial & Confidence triage & Yes & 1.000 & 19.2\% & 6.5\% & 0.500 \\
suffix1\_initial & Fixed one view & No & 1.000 & 0.0\% & 6.3\% & 0.370 \\
suffix1\_initial & Fixed two views & No & 2.000 & 0.0\% & 0.0\% & 0.348 \\
\hline
\end{tabular}}
\end{table*}

The conservative policy has a validation miss-risk point estimate below 5\% in both branches and saves images relative to fixed two-view acquisition. However, relative to matched confidence triage it acquires more images and defers more eyes. At unit deferral cost, the acquisition-minus-triage cost differences are 0.333 (paired patient-bootstrap 95\% interval 0.192--0.474) and 0.615 (0.412--0.838) image-equivalents per eye for suffix2 and suffix1 initial, respectively. These cost differences accompany different error rates and do not establish dominance or clinical utility. At a relaxed 10\% calibration target, acquisition can lower point-estimate cost for sufficiently expensive deferral, but the evaluated cost-difference intervals include zero at deferral cost 10 in both branches. The present comparison therefore does not demonstrate a robust cost advantage over triage.

\subsection{Formal Risk-Control Audit}
Table~\ref{tab:riskcontrol} reports the exact calibration-bound audit. No policy in the searched grid certifies a 5\% or 10\% patient-level missed-positive risk at the 95\% confidence level. At a 15\% target with a single-policy exact bound, first-view confidence triage is certifiable in both suffix branches; with Bonferroni correction over the full grid, no policy is certified. At 80\% confidence, confidence triage reaches a 10\% bound, but this is too weak for a clinical safety claim. This is an important result: the available public sample supports a strong benchmark paper, but it does not support claiming that the current adaptive acquisition model is risk-certified for deployment.

\begin{table*}[t]
\centering
\caption{Exact finite-sample calibration risk-control audit. Risk unit is patient-level missed referable DR among automatically decided positive-patient cases. UCB denotes one-sided Clopper--Pearson upper bound. Bonferroni correction is over 6,188 candidate policies per branch.}
\label{tab:riskcontrol}
\resizebox{\textwidth}{!}{\begin{tabular}{llllllll}
\hline
Branch & Confidence & Target & Correction & Certified method & Cal. UCB & Val. miss risk & Val. coverage \\
\hline
suffix2\_initial & 95\% & 15\% & single-policy & confidence triage & 0.133 & 0.118 & 0.923 \\
suffix1\_initial & 95\% & 15\% & single-policy & confidence triage & 0.146 & 0.111 & 0.821 \\
suffix2\_initial & 80\% & 10\% & single-policy & confidence triage & 0.074 & 0.118 & 0.923 \\
suffix1\_initial & 80\% & 10\% & single-policy & confidence triage & 0.081 & 0.111 & 0.821 \\
both branches & 95\% & 5--10\% & single-policy & none & -- & -- & -- \\
both branches & 95\% & 5--15\% & Bonferroni grid & none & -- & -- & -- \\
\hline
\end{tabular}}
\end{table*}

\subsection{Online Certification Attempt}
After the initial product build, we located the official DeepDRiD online challenge labels in the local v1.1 folder and converted them into a sealed certification-style cohort. This cohort contains 100 patients, 200 complete paired eyes, 400 images, 82 referable eyes, and 48 positive patients. The locked clinical artifact was evaluated without threshold tuning. It decided 47 positive patients and missed 7, giving a patient-level missed-risk point estimate of 14.9\% and a one-sided 95\% upper bound of 26.2\%. A post-hoc threshold sweep found zero-miss modes, but the best such mode decides only 48 positive patients, yielding a 95\% upper bound of 6.05\%, above the 5\% target. Therefore, the online cohort narrows the evidence gap but does not remove it: an autonomous 95\%/5\% certificate still needs at least 11 additional zero-miss positive-patient cases and, prospectively, a frozen operating rule tested on an untouched cohort.

\begin{table*}[t]
\centering
\caption{DeepDRiD online certification attempt with the locked clinical artifact. No threshold tuning was performed for the locked-artifact row. The post-hoc zero-miss row is diagnostic only and is not a prospective certificate.}
\label{tab:onlinecert}
\resizebox{\textwidth}{!}{\begin{tabular}{lllllll}
\hline
Mode & Positive patients & Decided positive patients & Misses & Miss risk & 95\% UCB & Decision \\
\hline
Locked artifact & 48 & 47 & 7 & 14.9\% & 26.2\% & no autonomous certificate \\
Best post-hoc zero-miss & 48 & 48 & 0 & 0.0\% & 6.05\% & still below 59-case requirement \\
Required zero-miss certificate & -- & 59 & 0 & 0.0\% & $\leq$5.0\% & target condition \\
\hline
\end{tabular}}
\end{table*}

\subsection{Frontier and Mechanism Tests}
The expanded benchmark sweeps calibration risk targets $\{0.01,0.02,0.05,0.10,0.15\}$ and minimum specificities $\{0.30,0.40,0.50,0.60\}$ for acquisition, confidence triage, fixed one-view, and fixed two-view baselines. The strict 1\% risk and 0.50 specificity setting is shown in Table~\ref{tab:policies}. Across this setting, acquisition reduces the observed conditional missed-risk point estimate from 6.3--6.5\% for confidence triage to 0.0--3.6\%, but coverage falls from 80.8--92.3\% to 60.3--79.5\% and mean imaging rises from 1.000 to 1.205--1.410 images per eye. Thus the result is not a simple win; it is a lower-risk, lower-coverage operating point.

\begin{figure*}[t]
\centering
\setlength{\unitlength}{1mm}
\begin{picture}(150,62)
\put(8,8){\line(1,0){132}}
\put(8,8){\line(0,1){48}}
\put(6,6){\scriptsize 0}
\put(24,3){\scriptsize 1.0}
\put(51,3){\scriptsize 1.2}
\put(78,3){\scriptsize 1.4}
\put(132,3){\scriptsize 2.0}
\put(52,0){\scriptsize mean images per eye}
\put(0,23){\scriptsize risk}
\put(8,22){\line(1,0){132}}
\put(118,24){\scriptsize 5\% reference}
\put(28,24){\circle*{3.5}}
\put(34,24){\scriptsize acquisition}
\put(18,27){\circle*{3.5}}
\put(24,27){\scriptsize triage}
\put(132,8){\circle*{3.5}}
\put(111,10){\scriptsize fixed two}
\put(18,26){\circle{5}}
\put(29,27){\circle{5}}
\put(38,26){\circle{5}}
\put(57,18){\circle{5}}
\put(87,8){\circle{5}}
\put(132,8){\circle{5}}
\put(18,33){\circle{5}}
\put(18,48){\circle{5}}
\end{picture}
\caption{Risk--image frontier schematic rendered from the validation operating-point table. Acquisition moves to lower missed-risk regions in parts of the frontier, while fixed two-view imaging reaches zero observed missed risk at two images per eye. The exact generated plot is saved as \texttt{plots/flagship\_risk\_image\_frontier.png}.}
\label{fig:frontier}
\end{figure*}

Cost-frontier bootstrapping reinforces the same conclusion. At strict calibration targets, acquisition is more expensive than confidence triage for $\lambda\in\{0,1,2,5,10\}$. At relaxed 10\% calibration risk, acquisition occasionally has a lower point-estimate cost when deferral is expensive, for example suffix2 at $\lambda=10$ gives $\Delta C=-0.051$, but the paired patient-bootstrap interval $[-0.321,0.103]$ includes zero. The current data therefore support acquisition as a plausible frontier component, not as a proven dominant policy.

\begin{figure*}[t]
\centering
\begin{tabular}{lllll}
\toprule
Branch and target & $\lambda=0$ & $\lambda=2$ & $\lambda=5$ & $\lambda=10$ \\
\midrule
suffix2, 1\% risk & 0.205 & 0.462 & 0.846 & 1.487 \\
suffix1, 1\% risk & 0.410 & 0.821 & 1.436 & 2.462 \\
suffix2, 10\% risk & 0.077 & 0.026 & 0.000 & -0.051 \\
suffix1, 10\% risk & 0.090 & 0.026 & -0.090 & -0.282 \\
\bottomrule
\end{tabular}
\caption{Operational cost sensitivity for acquisition relative to first-view confidence triage; negative values favor acquisition. Strict 1\% calibration-risk policies are costlier, while relaxed 10\% policies can approach or cross break-even when deferral is assigned high operational cost. The generated uncertainty plot is saved as \texttt{plots/flagship\_cost\_sensitivity.png}.}
\label{fig:cost}
\end{figure*}

\begin{table*}[t]
\centering
\caption{Mechanism diagnostics on validation. Negative Brier delta favors the paired view. The strongest reproducible mechanism is improved scoring in moderate DR or high score-discordance groups, not uniform improvement across all eyes.}
\label{tab:mechanism}
\resizebox{\textwidth}{!}{\begin{tabular}{lllllll}
\hline
Branch & Group & N & Ref. rate & Score delta & Brier delta & Helped / hurt \\
\hline
suffix2\_initial & grade 2 & 16 & 1.000 & 0.226 & -0.233 & 0.750 / 0.250 \\
suffix2\_initial & highest discordance quartile & 19 & 0.579 & 0.207 & -0.179 & 0.632 / 0.368 \\
suffix2\_initial & most uncertain quartile & 19 & 0.474 & 0.103 & -0.061 & 0.684 / 0.316 \\
suffix1\_initial & grade 1 & 9 & 0.000 & -0.049 & -0.098 & 0.556 / 0.444 \\
suffix1\_initial & uncertain quartile 2 & 19 & 0.316 & -0.062 & -0.080 & 0.579 / 0.421 \\
suffix1\_initial & grade 2 & 16 & 1.000 & -0.076 & 0.035 & 0.438 / 0.563 \\
\hline
\end{tabular}}
\end{table*}

\begin{figure*}[t]
\centering
\begin{tabular}{lll}
\toprule
Branch and subgroup & Brier delta & Visual magnitude \\
\midrule
suffix2 grade 2 & -0.233 & \rule{46mm}{2.5mm} \\
suffix2 high discordance & -0.179 & \rule{36mm}{2.5mm} \\
suffix2 most uncertain & -0.061 & \rule{12mm}{2.5mm} \\
suffix1 grade 1 & -0.098 & \rule{20mm}{2.5mm} \\
suffix1 high discordance & -0.038 & \rule{8mm}{2.5mm} \\
suffix1 grade 2 & 0.035 & worsens \\
\bottomrule
\end{tabular}
\caption{Mechanism diagnostics by clinically relevant subgroup. Negative Brier delta means the paired view improves score quality. The strongest gain occurs in suffix2-initial grade-2 eyes and high-discordance eyes, showing that adaptive acquisition value is concentrated rather than uniform. The generated PNG is saved as \texttt{plots/flagship\_mechanism\_brier.png}.}
\label{fig:mechanism}
\end{figure*}

The mechanism analysis also explains why mapping sensitivity matters. In the suffix2-initial branch, the second view strongly improves grade-2 eyes: mean score increases by 0.226 and Brier score improves by 0.233. In the suffix1-initial branch, grade-2 eyes worsen under the paired score by 0.035 Brier points, while nonreferable grade-1 eyes improve. A manuscript that names one suffix as a clinical field without official confirmation would therefore overstate the evidence. The defensible claim is mapping-robust paired-view benchmarking, not anatomy-specific field selection.

\subsection{Clinical Workflow Software Demonstration}
Table~\ref{tab:software} reports the software-side product artifacts generated from the same frozen evidence package. The locked artifact was trained on 792 DeepDRiD eyes with paired macula and optic-disc features. The image case runner was exercised on an actual paired DeepDRiD image case in two modes: frozen-cache inference and live ResNet-50 encoder inference with project-local model weights. The local browser station was also tested by uploading the same paired retinal images through its web form and by calling \texttt{POST /api/analyze}. It stored the uploaded files, wrote SHA-256 hashes, produced \texttt{result.json}, served \texttt{report.html}, and returned structured JSON containing referable probability, five-class grade prediction, clinician-facing impression, and recommended next step. The demo case stopped after the first view with a referable-DR score of 0.0053 and predicted grade 0, but the output still required manual review and explicitly disabled autonomous diagnosis. This is the intended product boundary: the software workflow is complete enough to operate on cases, audit inputs, and report decisions, while the model refuses unsupervised clinical action until risk certification is achieved.

\begin{table*}[t]
\centering
\caption{Clinical-workflow software artifacts and smoke evidence. These are product-readiness checks, not claims of autonomous clinical safety.}
\label{tab:software}
\resizebox{\textwidth}{!}{\begin{tabular}{llll}
\hline
Component & Artifact & Verification & Status \\
\hline
Locked model & \texttt{models/active\_retina\_rc\_clinical\_artifact.joblib} & trained on 792 paired eyes & ready for workflow testing \\
Metadata & \texttt{models/active\_retina\_rc\_clinical\_artifact.metadata.json} & manifest and feature-cache SHA-256 hashes recorded & reproducible \\
Gateway & \texttt{src/active\_retina/clinical\_gateway.py} & fail-closed tests passed & manual review required \\
Image case runner & \texttt{src/active\_retina/clinical\_case\_runner.py} & real paired-image demo executed offline & JSON and HTML written \\
Live encoder & \texttt{results/clinical\_case\_live\_encoder\_output.json} & project-local pretrained weights cached and used & new-image path verified \\
Browser station & \texttt{src/active\_retina/clinical\_station.py} & web upload generated \texttt{clinical\_station\_cases/WEB-DEMO-001} & user-facing path verified \\
Analysis API & \texttt{docs/api\_contract.md} & \texttt{POST /api/analyze} generated \texttt{results/api\_demo\_response.json} & hardware integration path verified \\
Demo report & \texttt{clinical\_station\_cases/WEB-DEMO-001/report.html} & clinician-readable report served locally & manual review banner present \\
Enclosure runbook & \texttt{docs/hardware\_enclosure\_clinical\_station.md} & operator workflow and enclosure constraints specified & product package defined \\
Autonomous candidate & \texttt{models/active\_retina\_rc\_autonomous\_candidate.joblib} & model family and thresholds frozen & awaits untouched certification \\
Certification runner & \texttt{src/active\_retina/certify\_autonomous\_candidate.py} & one-shot GO/NO-GO certificate emitted & computation path ready \\
Model card & \texttt{docs/autonomous\_candidate\_model\_card.md} & intended use and boundary specified & reproducible candidate package \\
\hline
\end{tabular}}
\end{table*}

\section{Discussion}
This study is strongest when read as a benchmark and falsification framework rather than as a completed deployment claim. The positive result is that paired views add measurable predictive information and that selective acquisition can move the operating point toward lower observed missed referable DR. The negative result is equally important: on the available development validation set, acquisition does not dominate a strong first-view confidence triage baseline once deferral, coverage, and uncertainty are counted. The DRTiD external run strengthens this conclusion by showing that a frozen DeepDRiD-derived autonomous candidate can fail sharply under paired-view domain shift even when the external cohort is large enough for finite-sample certification.

This is closer to a flagship scientific argument than a simple policy table because it exposes the central mechanism and the central failure mode. The mechanism is concentrated: the second view helps many moderate-disease and high-discordance eyes, but not every grade or mapping branch. The failure mode is also clear: if uncertain first-view cases can simply be deferred, then acquisition must justify its extra image burden by improving automatic resolution or risk in a way that survives patient-cluster uncertainty. The present data show plausibility, not proof.

\section{Publication Decision}
The current evidence supports a \emph{GO} decision for a transparent arXiv computational-system manuscript and for a bounded IEEE software/methods submission: the study has a reproducible clinical decision formulation, real paired-view predictions, five-fold patient-level cross-validation, inner calibration, binary and five-class endpoints, bootstrap uncertainty, mapping sensitivity, strong first-view triage baselines, generated figures, mechanism diagnostics, a machine-readable evidence audit, a finite-sample calibration-bound analysis, an online certification attempt, a locked case-running software layer, a browser station, a frozen autonomous-candidate artifact, and a one-shot certification runner. The contribution is not yet a certified autonomous clinical deployment. The contribution is that adaptive retinal imaging can now be evaluated and operated as a falsifiable risk--coverage--cost computational workflow with a frozen candidate ready for untouched certification.

The IEEE-journal decision is \emph{NO-GO} for an autonomous clinical claim and \emph{GO} for a bounded computational-system claim. The paper can now be framed around a locked model, locked thresholds, locked certification runner, public paired-view evidence, finite-sample risk arithmetic, and transparent failure modes. The DRTiD run shows that the frozen autonomous candidate is not externally certified. The correct claim is therefore a clinician-reviewed acquisition and analysis system plus a falsifiable certification framework, not autonomous clinical superiority or non-inferiority.

\section{Limitations}
The DeepDRiD suffix-to-anatomy mapping remains unofficial in the public files reviewed. The historical test split is exposed and excluded from confirmatory claims. The formal risk audit shows that the original calibration sample cannot certify a 5\% missed-DR risk target. The exact feasibility audit shows why: with zero misses, 59 automatically decided positive-patient calibration cases are needed for a one-sided 95\% upper bound below 5\%. The expanded development pool clears this arithmetic but is not untouched. DeepDRiD has two paired views, so all claims are paired-view workflow claims. The clinical software layer is workflow-ready but deliberately nonautonomous: it prevents unsupervised diagnostic action until the frozen candidate passes a prespecified untouched certification run.

\section{Conclusion}
ActiveRetina-RC provides a disciplined public-data route toward adaptive retinal imaging without relying on unavailable datasets. The GO result is now a full clinician-reviewed computational system: paired-view acquisition is judged by a risk--coverage--cost frontier, public DeepDRiD supports binary referable DR, five-class grading, calibrated acquisition policies, user-facing workflow software, and a frozen autonomous-candidate artifact. The candidate clears the finite-sample arithmetic in expanded development stress testing but fails external DRTiD certification. A final autonomous diagnosis claim therefore requires a stronger model and a new untouched certification run, with mBRSET remaining the most promising next patient-level route.


\section{Statistical Interpretation}
The current experiments impose empirical calibration risk and specificity constraints but no prespecified minimum coverage. This omission permits deferral to change the risk denominator, which is why risk, specificity, coverage, images, and deferral must be reported together. The finite-sample audit further shows that empirical zero-miss calibration behavior is not enough: with the current positive-patient calibration count, exact binomial bounds do not certify a 5\% risk target.

Let $I$ be images acquired, $D$ indicate deferral, and $Y$ denote referable DR. We compare $C_\lambda=\mathbb{E}[I]+\lambda\mathbb{E}[D]$, with $\lambda$ an explicitly hypothetical deferral cost in image-equivalents. For an acquisition policy $a$ and triage policy $t$, the break-even value is $(\mathbb{E}[I_a]-\mathbb{E}[I_t])/(\mathbb{E}[D_t]-\mathbb{E}[D_a])$ when the denominator is positive. This algebraic comparison is meaningful only together with error, specificity, and coverage reporting. A clinical objective would additionally account for false-positive referrals and missed disease; those costs were not measured here.

The validation cohort contains 78 eyes and 32 referable eyes. The conservative policies miss 1/28 decided referable eyes (3.6\%) in the suffix2-initial branch and 0/27 decided referable eyes (0.0\%) in the suffix1-initial branch. Specificity among decided eyes is 0.559 and 0.550, respectively. These point estimates meet the prespecified development gate, but they do not establish safety, equivalence, or noninferiority because the confidence intervals are wide and threshold selection is empirical. The more aggressive zero-deferral ablation misses 2/32 referable eyes (6.25\%) in both branches and is therefore not the primary risk-controlled policy.

Uncertainty is estimated by 1,000 paired patient-cluster bootstrap replicates with seed 20261003. Each sampled patient contributes all available eyes; the same resample evaluates acquisition and triage, and thresholds remain fixed. Percentile intervals describe development-sample variability conditional on the selected model and thresholds. They exclude model-training and threshold-selection variability and are not risk certificates. Full numeric intervals and operational cost sensitivity are provided in the reproducibility outputs.

There are 39 validation patients. For the conservative policy, the 95\% percentile interval for missed referable-eye risk is [0, 0.115] with suffix2 initial and [0, 0] with suffix1 initial. The suffix1 interval is zero because no misses occur in the bootstrap resamples of the observed validation set; it should not be read as a population guarantee. At the target-risk operating point, image savings remain 29.5--39.7\% versus always-two-view imaging. These intervals support candid reporting of uncertainty but do not establish equivalence.

\section{Related Work and Contribution Boundary}
Deep learning for diabetic retinopathy screening has been validated at large scale on retinal fundus photographs, including work on referable DR detection in EyePACS and Messidor settings\cite{gulshan2016} and multiethnic screening cohorts\cite{ting2017}. DeepDRiD extends this line with a public challenge dataset for DR grading, image quality estimation, and ultra-widefield grading\cite{deepdrid}. Two-field and multi-view retinal learning motivate studying whether multiple retinal views provide useful evidence beyond a single fundus image\cite{drtid}.

The decision rule studied here is related to selective prediction, in which a model may abstain from automatic classification to reduce risk at lower coverage\cite{selective,selectivenet}. It is also adjacent to active learning and active acquisition methods that use uncertainty or expected value to decide what information should be obtained next\cite{gal2017}. The distinction in this pilot study is that the acquisition action corresponds to an additional retinal image rather than an additional label. Our contribution is therefore not a new thresholding or selective-classification algorithm. It is an empirical formulation of physical second-view acquisition as part of a stop/acquire/defer screening workflow.

\section{Reproducibility}
The evidence package records source-prediction SHA-256 hashes, policy thresholds, mapping branches, bootstrap seed, and numeric outputs. The flagship evidence layer writes \texttt{flagship\_frontier.csv}, \texttt{flagship\_cost\_frontier.csv}, \texttt{flagship\_mechanism.csv}, \texttt{flagship\_oracle\_bounds.csv}, \texttt{flagship\_decision.json}, and generated figures for paired-view value, risk--image frontier, cost sensitivity, and mechanism diagnostics. The formal risk-control audit writes \texttt{formal\_risk\_control\_summary.json}, \texttt{formal\_risk\_control\_selected.csv}, and \texttt{formal\_risk\_control\_candidates.csv}. The feasibility audit writes \texttt{risk\_certificate\_feasibility.json}. The online certification attempt writes \texttt{deepdrid\_online\_certification\_manifest.csv}, \texttt{deepdrid\_online\_certification\_predictions.csv}, \texttt{deepdrid\_online\_certification\_summary.json}, and \texttt{certification\_mode\_search\_online.json}. The expanded development stress test writes \texttt{expanded\_cohort\_development\_summary.json} and \texttt{expanded\_cohort\_development\_table.csv}. The frozen candidate package writes \texttt{active\_retina\_rc\_autonomous\_candidate.joblib}, \texttt{active\_retina\_rc\_autonomous\_candidate.metadata.json}, \texttt{autonomous\_candidate\_model\_card.md}, \texttt{autonomous\_candidate\_online\_smoke\_predictions.csv}, and \texttt{autonomous\_candidate\_online\_smoke\_summary.json}. The full model benchmark writes \texttt{deepdrid\_full\_cv\_predictions.csv}, \texttt{deepdrid\_full\_cv\_metrics.csv}, \texttt{deepdrid\_full\_cv\_policy.csv}, and \texttt{deepdrid\_full\_cv\_folds.csv}. The clinical workflow package writes \texttt{active\_retina\_rc\_clinical\_artifact.joblib}, \texttt{active\_retina\_rc\_clinical\_artifact.metadata.json}, \texttt{clinical\_gateway\_example\_output.json}, \texttt{clinical\_gateway\_audit.jsonl}, \texttt{clinical\_case\_demo\_output.json}, \texttt{clinical\_case\_live\_encoder\_output.json}, and station case records under \texttt{clinical\_station\_cases/}. The historical exposed test partition contributes no confirmatory results. The next upgrade is no longer another benchmark sweep; it is the untouched certification run for the frozen candidate.

\begin{thebibliography}{9}
\bibitem{gulshan2016} V. Gulshan et al., ``Development and Validation of a Deep Learning Algorithm for Detection of Diabetic Retinopathy in Retinal Fundus Photographs,'' \emph{JAMA}, vol. 316, no. 22, pp. 2402--2410, 2016.
\bibitem{ting2017} D. S. W. Ting et al., ``Development and Validation of a Deep Learning System for Diabetic Retinopathy and Related Eye Diseases Using Retinal Images From Multiethnic Populations With Diabetes,'' \emph{JAMA}, vol. 318, no. 22, pp. 2211--2223, 2017.
\bibitem{deepdrid} ``DeepDRiD: Diabetic Retinopathy---Grading and Image Quality Estimation Challenge,'' \emph{Patterns}, vol. 3, no. 6, 2022. Available: \url{https://pmc.ncbi.nlm.nih.gov/articles/PMC9214346/}.
\bibitem{drtid} ``Cross-Field Transformer for Diabetic Retinopathy Grading on Two-Field Fundus Images,'' 2022. Available: \url{https://arxiv.org/abs/2211.14552}.
\bibitem{mbrset} L. F. Nakayama et al., ``mBRSET, a Mobile Brazilian Retinal Dataset,'' \emph{PhysioNet}, version 1.0, 2024. doi: 10.13026/qxpd-1y65.
\bibitem{selective} Y. Geifman and R. El-Yaniv, ``Selective Classification for Deep Neural Networks,'' \emph{NeurIPS}, 2017. Available: \url{https://arxiv.org/abs/1705.08500}.
\bibitem{selectivenet} Y. Geifman and R. El-Yaniv, ``SelectiveNet: A Deep Neural Network with an Integrated Reject Option,'' \emph{ICML}, 2019. Available: \url{https://arxiv.org/abs/1901.09192}.
\bibitem{gal2017} Y. Gal, R. Islam, and Z. Ghahramani, ``Deep Bayesian Active Learning with Image Data,'' \emph{ICML}, pp. 1183--1192, 2017.
\end{thebibliography}

\end{document}
